% CHECKPOINT36: complete formalization of frozen research source
% actual_elliptic_pentagon_cubic35.md
% SHA256 ab9ccffba8f12ce224aab9fd3c812bfdf8b24d8117b61bf28c77490b291b196d.
\statement{A.121}{Proved}{The full profinite elliptic section and a pentagon-only marked cubic}
Write \(G=\widehat B_{1,3}\) for the full profinite completion of the
reduced elliptic braid group, with the original marking specified below.
Put
\[
 K=[Y,X]=\sigma\tau^2\sigma,\qquad L=\tau^{-1}\sigma^{-1}.
\]
For every profinite commutator word \(w\) satisfying the original
four-strand pentagon, the following literal identity holds in \(G\):
\[
 \Pi_w^3=1,\qquad
 \Pi_w=L\,w(YX^{-1}Y^{-1},K)\,Y\,w(X,K)^{-1}
                      w(\sigma^2,\tau^2).                    \tag{A.121.1}
\]
This assertion requires neither a hexagon for \(w\), a scalar multiplier,
nor an automorphism associated with \(w\). It applies to the actual
canonical word \(u\) of a negative involution of full
\(\widehat{\mathrm{GT}}\), by A.37. Its exact transported form under
A.118 is (A.121.13).

\begin{proof}
\emph{1. The primary full profinite section and the actual marking.}
We use the convention \([a,b]=aba^{-1}b^{-1}\).
The full profinite case of
\citet[Proposition~2.20]{A_Enriquez2014} gives a section of the elliptic
to ordinary GT semigroup map. For an ordinary GT pair \((\lambda,f)\)
its two words, with \(m=(\lambda-1)/2\), are
\[
 \begin{aligned}
 g_+(X,Y)&=f(X,K)X^\lambda f(X,K)^{-1},\\
 g_-(X,Y)&=K^m f(YX^{-1}Y^{-1},K)Yf(X,K)^{-1}.
 \end{aligned}                                                \tag{A.121.2}
\]
The section is unital and multiplicative, so it takes units to units
and involutions to involutions. Its action on the two-generator
abelianization is \(\operatorname{diag}(\lambda,1)\).
The cited paper fixes the second braid-coordinate generator and
conjugates the first. Interchanging the two generators and using
duality gives the project's normalized convention
\(x\mapsto x^\lambda,\ y\mapsto f^{-1}y^\lambda f\), with the
same word \(f\). The opposite-composition convention in the paper's
endomorphism factor is retained. We do not apply this section to
the still unproved GT pair \((d,u)\), or to the GT-zero lift of A.118.

The presentation used here is the \(n=3\) case of
\citet[Section~1.2, equation~(7)]{A_Enriquez2014}. Set
\(\sigma=\sigma_1,\ \tau=\sigma_2,\ X=X_1^+,\ Y=X_1^-\).
It gives
\[
 \begin{gathered}
 \sigma\tau\sigma=\tau\sigma\tau,\qquad
 [X,\tau]=[Y,\tau]=1,\\
 (\sigma X)^2=(X\sigma)^2,\qquad
 (\sigma^{-1}Y)^2=(Y\sigma^{-1})^2,\qquad
 [Y,X_2^{-1}]=\sigma^2,\\
 X_2=\sigma X\sigma,\quad X_3=\tau X_2\tau,\quad XX_2X_3=1,\\
 Y_2=\sigma^{-1}Y\sigma^{-1},\quad
 Y_3=\tau^{-1}Y_2\tau^{-1},\quad YY_2Y_3=1 .
 \end{gathered}                                                \tag{A.121.3}
\]
The same-sign \(X_i\) commute, as do the \(Y_i\), and
\(K=[Y,X]=\sigma\tau^2\sigma\). Put \(K_2=\tau\sigma^2\tau\).
The type-B braid homomorphism in the proof of Proposition~2.20 is
\[
 (B_4/Z(B_4))\times_{S_4}S_3\longrightarrow B_{1,3},\qquad
 \sigma_0^2\longmapsto X,\quad
 \sigma_1\longmapsto\sigma,\quad\sigma_2\longmapsto\tau.          \tag{A.121.4}
\]
Only this homomorphism is needed. No injectivity or completion
recognition is inferred. The homomorphism extends to full profinite
completions by their universal property.

\emph{2. The original pentagon, with every argument retained.}
In the original \(P_4\) marking write
\((a,b,c,d_{\mathrm{br}},e,f_{\mathrm{br}})
=(A_{12},A_{13},A_{23},A_{14},A_{24},A_{34})\).
Here \(a=\sigma_0^2,\ c=\sigma_1^2,\ f_{\mathrm{br}}=\sigma_2^2\),
and
\[
 \begin{gathered}
 ce=\sigma_1\sigma_2^2\sigma_1,\qquad
 bc=\sigma_1\sigma_0^2\sigma_1,\\
 ab=\sigma_0\sigma_1^2\sigma_0,\qquad
 ef_{\mathrm{br}}=\sigma_2\sigma_1^2\sigma_2 .
 \end{gathered}
\]
Thus the given pentagon is literally
\[
 w(a,ce)w(bc,f_{\mathrm{br}})
 =w(c,f_{\mathrm{br}})w(ab,ef_{\mathrm{br}})w(a,c).
\]
In the middle factor \(\sigma_1^2\) commutes with both arguments and
\((\sigma_1\sigma_0^2)^2
 =(\sigma_0\sigma_1^2\sigma_0)\sigma_1^2\).
A commutator word is unchanged by these commuting decorations:
this follows first for ordinary words and then by continuity.
Applying (A.121.4), and using \(X_2X=X_3^{-1}\), gives exactly
\[
 w(\sigma^2,\tau^2)w(X_3^{-1},K_2)w(X,\sigma^2)
       =w(X,K)w(X_2,\tau^2).                                  \tag{A.121.5}
\]

\emph{3. Ordered conjugation to a three-factor norm.}
The marked relations give
\[
 L(YX^{-1}Y^{-1})L^{-1}=X_3^{-1},\qquad LKL^{-1}=K_2 .
\]
Writing \(E=w(X_3^{-1},K_2)\), we have
\(\Pi_w=E L Yw(X,K)^{-1}w(\sigma^2,\tau^2)\).
Conjugate first by \(E^{-1}\), then by \(L^{-1}\), and use
(A.121.5). The result is
\[
 T_w=Y A_wL,\qquad
 A_w=w(X_2,\tau^2)w(X,\sigma^2)^{-1}.                          \tag{A.121.6}
\]
Set \(B=\tau\sigma,\ C=\sigma\tau\). Since
\(LY=Y_3B,\ BLY=Y_2C,\ CL=1\), direct multiplication gives
\[
 T_w^3=Y A_wY_3\,\operatorname{Ad}(B)(A_w)
                       Y_2\,\operatorname{Ad}(C)(A_w).        \tag{A.121.7}
\]
No central full twist or cyclic power has been removed.

\emph{4. Complete cancellation in the marked group.}
Define
\[
 \begin{aligned}
 B_w&=w(X_2\sigma^{-2},\tau^2),\\
 D_w&=w(\tau\sigma^2\tau^{-1}X,\sigma^2),\\
 E_w&=w(X_3\tau^{-1}\sigma^{-2}\tau^{-1},
                                     \tau\sigma^2\tau^{-1}),\\
 F_w&=w(X_3\tau^{-1}\sigma^{-2}\tau,
                                     \tau\sigma^2\tau^{-1}).
 \end{aligned}                                                \tag{A.121.8}
\]
The required literal argument transformations from (A.121.3) are
\[
 \begin{gathered}
 YX_2=(X_2\sigma^{-2})Y,\qquad [Y,\tau^2]=1,\\
 XY_3=Y_3(\tau\sigma^2\tau^{-1}X),\qquad
 [\sigma^2,Y_3]=1,\qquad YY_3=Y_2^{-1},\\
 \operatorname{Ad}(B)(X_2)=\tau\sigma^2\tau^{-1}X,\qquad
 \operatorname{Ad}(B)(\tau^2)=\sigma^2,\\
 \operatorname{Ad}(B)(X)=X_3\tau^{-1}\sigma^{-2}\tau^{-1},\qquad
 \operatorname{Ad}(B)(\sigma^2)=\tau\sigma^2\tau^{-1}.
 \end{gathered}                                                \tag{A.121.9}
\]
The remaining two images under \(\operatorname{Ad}(C)\) are
\[
 \begin{gathered}
 \operatorname{Ad}(C)(X_2)=X_3\tau^{-1}\sigma^{-2}\tau,\qquad
 \operatorname{Ad}(C)(\tau^2)=\tau\sigma^2\tau^{-1},\\
 \operatorname{Ad}(C)(X)=X_2\sigma^{-2},\qquad
 \operatorname{Ad}(C)(\sigma^2)=\tau^2 .
 \end{gathered}
\]
Substitution, with the indicated order, gives
\[
 \begin{aligned}
 YA_wY_3&=B_wY_2^{-1}D_w^{-1},\\
 \operatorname{Ad}(B)(A_w)&=D_wE_w^{-1},\\
 \operatorname{Ad}(C)(A_w)&=F_wB_w^{-1}.
 \end{aligned}                                                \tag{A.121.10}
\]
Finally,
\[
 \begin{gathered}
 Y_2^{-1}(X_3\tau^{-1}\sigma^{-2}\tau^{-1})Y_2
              =X_3\tau^{-1}\sigma^{-2}\tau,\\
 [Y_2,\tau\sigma^2\tau^{-1}]=1,
 \end{gathered}
\]
so \(Y_2^{-1}E_wY_2=F_w\). Equations (A.121.7) and
(A.121.10) now give
\[
 T_w^3=B_wY_2^{-1}E_w^{-1}Y_2F_wB_w^{-1}=1 .
\]
Conjugacy of \(T_w\) and \(\Pi_w\) proves (A.121.1).
These are discrete group identities followed by continuous evaluation
of the given profinite word; no finite quotient census is used.

\emph{5. Canonical and transported specializations.}
For a negative involution of full \(\widehat{\mathrm{GT}}\), A.37
supplies the pentagon for its actual canonical \(u\). For any
\(m\in\widehat{\mathbb Z}\), the candidate second word
\[
 g_-[w,m]=K^m w(YX^{-1}Y^{-1},K)Yw(X,K)^{-1}
\]
satisfies
\[
 \bigl(LK^{-m}g_-[w,m]w(\sigma^2,\tau^2)\bigr)^3=1,            \tag{A.121.11}
\]
because its adjacent \(K\)-powers cancel. In particular \(m=k\)
is allowed for the actual canonical \(d=1+2k\).
The other defining elliptic relations are not inferred: the primary
proof of the positive cubic uses the usual GT braid action, and its
commutator-relation proof uses the hexagon.

For the actual \(\Gamma=(d,f)\) and \(U'\) of A.118, define for
every pair \((p,q)\) in a profinite target
\[
 \mathcal U_{p,q}
 =f(p,q)\,
 U'\bigl(p^d,f(p,q)^{-1}q^df(p,q)\bigr).                        \tag{A.121.12}
\]
The equality \(u=f\Gamma(U')\) gives
\(\mathcal U_{p,q}=u(p,q)\) with this exact order. Hence
\[
 \left(
 L\,\mathcal U_{YX^{-1}Y^{-1},K}\,Y\,
       \mathcal U_{X,K}^{-1}\mathcal U_{\sigma^2,\tau^2}
 \right)^3=1.                                                 \tag{A.121.13}
\]
This is a pairwise transported identity. It neither assigns an
elliptic action to \(\Gamma\) nor deletes its \(f\)-factors.
It does not assert an ordinary cubic for \(U'\), the boundary equation
of A.122, or the original coherent-frame or CM conclusion.
\end{proof}
